\subsection{方阵}

定义 4. 行与列的指标集相同的矩阵称为方阵。

具有 n 行 n 列的方阵称为 n 阶矩阵。

注记。应当指出，行与列的指标集具有相同基数但并不相同的矩阵，绝不可视为方阵；特别地，这样的两个矩阵在环上的乘积是没有定义的。

显然，A 上的、以有限集作为行与列指标集的方阵的加法与乘法，由于公式 (7)、(8) 与 (9)（第 3 号），在这些矩阵的集合上定义了一个环结构；矩阵 \((\delta_{ij})\) （其中 \(\delta_{ij}\) 是 Kronecker 指标（对 \(i \in I\) , \(j \in I\) ））是该环的单位元，记作 \(I_n\) 或 \(1_n\) （当 I 有 \(n\) 个元素时）。当 \(\mathbf{I} = [1, n]\) 时，如此定义的矩阵环简记为 \(\mathbf{M}_n(\mathbf{A})\) ； \(\mathbf{M}_n(\mathbf{A})\) 的可逆元素群记作 \(\mathbf{GL}_n(\mathbf{A})\) 或 \(\mathbf{GL}(n, \mathbf{A})\) 。

A 上的 n 阶方阵 \(U = (a_{ij})\) 为右（相应地左）可逆的必要且充分条件是：对 A 的元素的每个组 \((b_i)_{1 \leq i \leq n}\) ，由 n 个未知量的 n 个方程构成的方程组

\[
\sum_ {j = 1} ^ {n} a _ {i j} x _ {j} = b _ {i} \quad (1 \leqslant i \leqslant n)
\]

\[
\left(\text {resp.} \sum_ {j = 1} ^ {n} x _ {j} a _ {j i} = b _ {i}\right)
\]

都在 A 中有一个解 \((x_{i})\) 。

设 I 是有限指标集，A 是环，E 是具有基 \((e_{i})_{i\in I}\) 的右（相应地左）A-模。对 E 的每个自同态 u，矩阵 \(M(u)\) （u 关于两个都等同于 \((e_{i})\) 的基的）是一个方阵；更简单地说，它称为 u 关于基 \((e_{i})\) 的矩阵。

设 \(\mathbf{I} = [1, n]\) 。映射 \(u \mapsto M(u)\) （相应地 \(u \mapsto {}^t M(u)\) ）是环 \(\operatorname{End}_{\mathbf{A}}(\mathbf{E})\) 到 \(\mathbf{M}_n(\mathbf{A})\) （相应地到 \(\mathbf{M}_n(\mathbf{A})\) 的反环上，这由公式 (18 D)（相应地 (18 G)）得出）（第 4 号）的一个同构。环 \(\mathbf{M}_{n}(\mathbf{A})\) 的可逆元素称为可逆矩阵，它们在映射 \(u \mapsto M(u)\) （相应地 \(u \mapsto {}^{t}M(u)\) ）之下对应于 E 的自同构；因此群 \(\mathbf{GL}(n, \mathbf{A})\) 典范地等同于群 \(\mathbf{GL}(\mathbf{A}_{d}^{n})\) 。

若 \(u\) 是 \(\mathbf{E}\) 的一个自同构，则它的逆步 \(\check{u}\) 是左（相应地右）A-模 \(\mathbf{E}^*\) 的一个自同构，使得 \(\check{u} = (^{t}u)^{-1} = ^{t}(u^{-1})\) （§ 2，第 5 号，定义 6）；若 \(M(\check{u})\) 是 \(\check{u}\) 关于对偶基 \((e_{i}^{*})\) 的矩阵，则由命题 3（第 4 号）可得

\[
M (\check {u}) = (^ {t} M (u)) ^ {- 1} = ^ {t} M (u ^ {- 1}).\tag{29}
\]

因此，对每个可逆矩阵 X，有 \({}^{t}(X^{-1}) = ({}^{t}X)^{-1}\) ；这个矩阵也记作 \({}^{t}X^{-1}\) ，称为矩阵 X 的逆步矩阵。

设 \(\sigma\) 是环 A 的一个自同构；对每个半线性映射 \(u: \mathbf{E} \to \mathbf{E}\) （相对于 \(\sigma\) ），矩阵 \(M(u)\) （该映射关于 E 的基 \((e_i)\) 的）也是一个方阵。由 (27 D)（第 6 号）立即可得，若 \(u\) 是双射，则

\[
M (u ^ {- 1}) = (\sigma^ {- 1} (M (u))) ^ {- 1}.
\]

设 E 是一个 A-模，它是子模的有限族 \((\mathbf{E}_{i})_{i\in I}\) 的直和；对 E 的每个自同态 u，矩阵 \(M(u)=(u_{kt})\) （u 关于 E 的两个都等同于 \((\mathbf{E}_{i})\) 的分解（第 5 号）的）是一个线性映射的方阵。要使 \(u(\mathbf{E}_{i})\subset\mathbf{E}_{i}\) 对所有 \(i\in I\) 成立，必要且充分条件是 \(u_{kt}=0\) （对 \(k\neq i\) ）。当 \(\mathbf{I} = [1, n]\) 时，关系

\[
u (\mathrm{E} _ {i}) \subset \mathrm{E} _ {i} + \mathrm{E} _ {i + 1} + \dots + \mathrm{E} _ {n} \quad (1 \leqslant i \leqslant n)
\]

等价于关系 \(u_{ki} = 0\) （对 \(k < i\) ）。

方阵的例子。I. 对角矩阵。在一个方阵中，

\[
M = \left(m _ {\iota \kappa}\right) _ {(\iota, \kappa) \in I \times I},
\]

两个指标相等的元素称为对角元素，族 \((m_{\iota})_{\iota \in I}\) 称为 \(M\) 的对角线；环上的方阵 \(M = (m_{\iota \kappa})\) ，若其对角元素以外的元素全为零，则称为对角矩阵。对环 A 的元素的每个族 \((a_{\iota})_{\iota \in I}\) ，对角矩阵 \((m_{\iota \kappa})\) （满足 \(m_{\iota} = a_{\iota}\) （对所有 \(\iota \in I\) ））记作 \(\operatorname{diag}(a_{\iota})_{\iota \in I}\) （或 \(\operatorname{diag}(a_1, a_2, \ldots, a_n)\) ，当 \(I = [1, n]\) 时）。在集合 \(\mathbf{M}_n(A)\) （A 上 \(n\) 阶方阵的）中，单位矩阵 \(I_n\) 是对角矩阵，并且它的每个倍数 \(aI_n = I_na\) （该矩阵被纯量 \(a\) 乘）也是（即称为标量矩阵的对角矩阵，其所有对角元素都等于 \(a\) ）。

对 A 的元素的每个族 \((d_i)_{1\leqslant i\leqslant n}\) 以及每个矩阵 \(X = (x_{ij})\) （ \((n,q)\) 型（相应地 \((p,n))\) ），A 上的），记 \(D = \mathrm{diag}(d_i)\) ，则

\[
\left\{ \begin{array}{l} D X = (d _ {i} x _ {i j}) \\ X D = (x _ {i j} d _ {j}). \end{array} \right.\tag{30}
\]

特别地，对两个 n 阶对角矩阵，有

\[
\begin{array}{c} \operatorname{diag} (a _ {i}) + \operatorname{diag} (b _ {i}) = \operatorname{diag} (a _ {i} + b _ {i}) \\ \operatorname{diag} (a _ {i}). \operatorname{diag} (b _ {i}) = \operatorname{diag} (a _ {i} b _ {i}). \end{array}\tag{31}
\]

因此对角矩阵构成 \(\mathbf{M}_{n}(\mathbf{A})\) 的一个子环，它同构于乘积环 \(A^{n}\) ；标量矩阵构成一个同构于 A 的子环。

\begin{enumerate}
\def\labelenumi{\Roman{enumi}.}
\setcounter{enumi}{1}
\tightlist
\item
  置换矩阵；单项矩阵。设 \(\pi\) 是有限集 I 的任意一个置换，并设 \((e_i)_{i\in I}\) 是 A-模 \(\mathbf{E} = \mathbf{A}_d^{\mathbf{I}}\) 的典范基；存在且只存在一个 E 的自同态 \(u_{\pi}\) ，使得对所有 \(i\in I\) , \(u_{\pi}(e_i) = e_{\pi(i)}\) （§ 1，第 11 号，命题 17 的推论 3）。对所有 \(i\in I\) ，指标为 \(i\) 的列（在矩阵 \(M(u_{\pi})\) 中，关于基 \((e_i)\) ）除指标为 \(\pi(i)\) 的行中那个等于 1 的元素外，其余元素全为零。按一种习用说法， \(M(u_{\pi})\) 称为置换 \(\pi\) 的矩阵。立即可见，对 I 的任意两个置换 \(\sigma, \tau\) ，有 \(u_{\sigma \tau} = u_{\sigma} \circ u_{\tau}\) ，且对恒等置换 \(\varepsilon, u_{\varepsilon}\) 为恒等映射；因此映射 \(\pi \mapsto M(u_{\pi})\) 是对称群 \(\mathfrak{S}_{\mathbf{I}}\) 到置换矩阵群上的一个同构。
\end{enumerate}

置换矩阵的每一行与每一列只含一个元素 \(\neq0\) 。非零环 A 上的、具有此性质的有限方阵 R 称为单项矩阵；设 \(r_{i}\) 是 R 的指标为 i 的列中那个唯一的元素 \(\neq0\) ，并设 \(\pi(i)\) 为该元素所在行的指标；显然 \(\pi\) 是指标集 I 的一个置换，且 \(R=M(u_{\pi})D\) ，其中 \(D=\mathrm{diag}(r_{i})\) 。

\begin{enumerate}
\def\labelenumi{\Roman{enumi}.}
\setcounter{enumi}{2}
\tightlist
\item
  三角矩阵。在环 A 上 n 阶方阵的环 \(\mathbf{M}_{n}(\mathbf{A})\) 中，任意矩阵 \((a_{ij})\) ，若 \(a_{ij}=0\) （对 i\textgreater j（相应地 i\textless j）），则称为上（相应地下）三角矩阵；也说这样的矩阵在其对角线以下（相应地以上）只有零。立即可证，上（相应地下）三角矩阵构成 \(\mathbf{M}_{n}(\mathbf{A}),\mathbf{S}\cap\mathbf{T}\) 的一个子环 S（相应地 T），后者显然就是对角矩阵环。
\end{enumerate}

S（相应地 T）中对角元素可逆的矩阵所成的集合 S'（相应地 T'）是一个矩阵乘法群，称为上（相应地下）全三角群，这由 § 1 第 11 号注记 5 直接得出。S（相应地 T）中对角元素全等于 1 的矩阵所成的集合 \(S_{1}\) （相应地 \(T_{1}\) ）是上述群的一个子群，称为上（相应地下）严格三角群，且 S'（相应地 T'）中对角线为 \((d_{i})\) 的每个矩阵 M 都可写成 \(M = DM_{1} = M_{1}'D\) ，其中 \(D = \mathrm{diag}(d_{i})\) ，而 \(M_{1}\) 与 \(M_{1}'\) 是属于 \(S_{1}\) （相应地 \(T_{1}\) ）的矩阵。

\begin{enumerate}
\def\labelenumi{\Roman{enumi}.}
\setcounter{enumi}{3}
\tightlist
\item
  矩阵的对角矩阵与三角矩阵。设 \((\mathbf{I}_k)_{1 \leqslant k \leqslant p}\) 是有限集 I 的一个划分；环 A 上以 I 为指标集的每个方阵都可以写成一个矩阵的方阵的形式，它对应于行的指标集与列的指标集的同一划分 \((I_{k})\) （第 5 号）
\end{enumerate}

\[
\left( \begin{array}{c c c c} X _ {1 1} & X _ {1 2} & \ldots & X _ {1 p} \\ X _ {2 1} & X _ {2 2} & \ldots & X _ {2 p} \\ \cdot & \cdot & \cdot & \cdot \\ X _ {p 1} & X _ {p 2} & \ldots & X _ {p p} \end{array} \right)\tag{32}
\]

其中每个 \(X_{kk}\) 是以 \(I_{k}\) 为行和列指标集的方阵。

在此记号下，若所有矩阵 \(X_{ij}\) （满足 \(i \neq j\) （相应地 i \textgreater{} j，相应地 i \textless{} j）的）均为零，则称 (32) 为矩阵的对角矩阵（相应地矩阵的上三角矩阵、相应地矩阵的下三角矩阵）。其矩阵为矩阵的对角矩阵、相应地三角矩阵的自同态 u 的解释，前面已通过考虑相应的线性映射的矩阵 \(M(u)\) 而看到过。对 I 的给定划分 \((\mathbf{I}_{k})\) ，矩阵的下三角（相应地上三角、对角）矩阵构成矩阵环 \(A^{I \times I}\) 的一些子环。特别地，相对于

划分 \((\mathbf{I}_k)\) 的矩阵的对角矩阵环同构于乘积 \(\prod_{k=1}^{p} \operatorname{End}_{\mathbf{A}}(\mathbf{E}_k)\) 。
% semantic-fragments: legacy-input-seam
\relax{}
